% TOP_PROBLEM: {"id": 9500003, "title": "Percolation dimension of planar Brownian trace", "category_id": 19, "category": {"id": 19, "name": "probability", "display_name": "Probability", "description": "Problems involving probability theory, stochastic processes, and random structures.", "slug": "probability", "order_index": 19, "created_at": "2026-07-31T15:26:25.671Z"}, "status": "open", "problem_number": "AMR-094-0003", "set_id": 13}
% FIRST_POSTED: 2026-10-01
% ENTRY_KIND: statement_only
% UnsolvedMath ID 9500003; source catalogue code AMR-094-0003
% !TeX program = lualatex
% Definition and sources only; this file is not an LLM solution attempt.
\documentclass[11pt,a4paper]{article}
\usepackage[margin=25mm]{geometry}
\usepackage{fontspec}
\setmainfont{lmroman10-regular.otf}[BoldFont=lmroman10-bold.otf,
 ItalicFont=lmroman10-italic.otf,BoldItalicFont=lmroman10-bolditalic.otf]
\usepackage{amsmath,amssymb,mathtools}
\usepackage{unicode-math}
\setmathfont{latinmodern-math.otf}
\AtBeginDocument{\let\setminus\smallsetminus}
\usepackage{xurl,hyperref}
\hypersetup{colorlinks=true,urlcolor=blue}
\setcounter{secnumdepth}{0}
\setlength{\parindent}{0pt}
\setlength{\parskip}{5pt}
\setlength{\emergencystretch}{3em}
\begin{document}
\section*{Percolation dimension of planar Brownian trace}\label{9500003}
{\small UnsolvedMath ID 9500003 · AMR-094-0003 · Probability · AMR Open Problem Lists}
\subsection{Definitions and mathematical statement}
Let $X=(X_t)_{0\le t\le1}$ be standard Brownian motion in $\mathbb R^2$, and write $T=X[0,1]$ for its trace. For a set $E\subseteq\mathbb R^2$, define its Hausdorff dimension by
\[
\dim_H E=\inf\{s\ge0:\mathcal H^s(E)=0\},
\]
where $\mathcal H^s$ is the Hausdorff measure obtained by taking infima of sums of $s$th powers of covering diameters and letting the permitted diameters tend to zero.

A nondegenerate Jordan arc is a subset homeomorphic to $[0,1]$ with two distinct endpoints. For any set $B$, its percolation dimension in the source's terminology is
\[
\operatorname{pdim}(B)=\inf\{\dim_H A:A\subseteq B,
\ A\text{ is a nondegenerate Jordan arc}\}.
\]
Take the infimum of an empty family to be $+\infty$. Is $\operatorname{pdim}(T)=1$ almost surely?

Every nondegenerate connected set has Hausdorff dimension at least one, so the substantive demand is the existence, for every $\varepsilon>0$, of some Jordan arc in the trace with dimension less than $1+\varepsilon$. Such arcs may vary with $\varepsilon$ and with the realization of Brownian motion. The conjecture does not require the infimum to be attained by an arc of dimension exactly one. Nor does it prescribe the endpoints or require an arc joining the endpoints of the complete Brownian path.
\subsection{Short English statement}
Does a planar Brownian trace contain Jordan arcs of Hausdorff dimension arbitrarily close to one?
\subsection{Sources}
\begin{itemize}
\item[S1] ULAM.AI and the UnsolvedMath Contributors, supplied problems.json, record 9500003 (AMR-094-0003), AMR Open Problem Lists. Catalogue identity and source transcription; CC BY 4.0.
\url{https://huggingface.co/datasets/ulamai/UnsolvedMath}.
\item[S2] Burdzy - My favorite open problems (2009), Problem 3. Original source indexed by AMR-094-0003.
\url{https://sites.math.washington.edu/~burdzy/open_mathjax.php}.
\end{itemize}
\end{document}
